% Circuito de sinal da Questão 2 (base-comum, PNP) com o transistor no modelo T, na região sombreada.
% Compilar com pdflatex; o PDF resultante substitui figuras/fig_q2_T.pdf.
\documentclass[11pt,border=4pt]{standalone}
% Estilo comum dos esquemáticos do relatório (circuitikz).
% Uso: \documentclass[11pt,border=4pt]{standalone} e, em seguida, % Estilo comum dos esquemáticos do relatório (circuitikz).
% Uso: \documentclass[11pt,border=4pt]{standalone} e, em seguida, % Estilo comum dos esquemáticos do relatório (circuitikz).
% Uso: \documentclass[11pt,border=4pt]{standalone} e, em seguida, \input{estilo_circuitos}.
\usepackage[T1]{fontenc}
\usepackage[utf8]{inputenc}
\usepackage[brazilian]{babel}
\usepackage{lmodern}
\usepackage{amsmath}
\usepackage[american]{circuitikz}
\usetikzlibrary{calc,arrows.meta}

% Cores das figuras originais
\definecolor{anot}{RGB}{0,135,200}      % anotações: valores calculados, setas, resistências vistas
\definecolor{sombra}{RGB}{220,242,253}  % região do transistor nos modelos T e pi
\definecolor{term}{RGB}{62,124,148}     % nomes de terminais destacados (B, C, E)

% Fios com 0,4 pt e corpo dos componentes com 0,8 pt, como nas figuras originais
\ctikzset{bipoles/thickness=2}
.
\usepackage[T1]{fontenc}
\usepackage[utf8]{inputenc}
\usepackage[brazilian]{babel}
\usepackage{lmodern}
\usepackage{amsmath}
\usepackage[american]{circuitikz}
\usetikzlibrary{calc,arrows.meta}

% Cores das figuras originais
\definecolor{anot}{RGB}{0,135,200}      % anotações: valores calculados, setas, resistências vistas
\definecolor{sombra}{RGB}{220,242,253}  % região do transistor nos modelos T e pi
\definecolor{term}{RGB}{62,124,148}     % nomes de terminais destacados (B, C, E)

% Fios com 0,4 pt e corpo dos componentes com 0,8 pt, como nas figuras originais
\ctikzset{bipoles/thickness=2}
.
\usepackage[T1]{fontenc}
\usepackage[utf8]{inputenc}
\usepackage[brazilian]{babel}
\usepackage{lmodern}
\usepackage{amsmath}
\usepackage[american]{circuitikz}
\usetikzlibrary{calc,arrows.meta}

% Cores das figuras originais
\definecolor{anot}{RGB}{0,135,200}      % anotações: valores calculados, setas, resistências vistas
\definecolor{sombra}{RGB}{220,242,253}  % região do transistor nos modelos T e pi
\definecolor{term}{RGB}{62,124,148}     % nomes de terminais destacados (B, C, E)

% Fios com 0,4 pt e corpo dos componentes com 0,8 pt, como nas figuras originais
\ctikzset{bipoles/thickness=2}

\begin{document}
\begin{circuitikz}
  % ---- posições (cm) ----
  \def\yC{5.75}    % fio do coletor
  \def\yB{2.45}    % nó interno da base
  \def\yS{3.75}    % centro da fonte controlada
  \def\yRo{5.20}   % altura da seta de R_o
  \def\xJ{2.85}    % nó de R_S, R_E e emissor
  \def\xE{3.85}    % terminais E e B
  \def\xBg{3.19}   % terra da base
  \def\xVe{4.24}   % anotação de v_e
  \def\xT{5.17}    % eixo do modelo T
  \def\xA{5.62}    % seta de i_e
  \def\xSl{4.35}   % região sombreada, borda esquerda
  \def\xSr{7.13}   % região sombreada, borda direita
  \def\xC{7.58}    % terminal C
  \def\xRC{8.68}   % R_C
  \def\xRo{9.70}   % seta de R_o
  \def\xRL{10.45}  % R_L
  \def\xVo{11.10}  % terminal de saída
  \def\xRin{2.02}  % seta de R_in

  % ---- transistor (modelo T) ----
  \fill[sombra] (\xSl,-0.38) rectangle (\xSr,\yC+0.33);

  % ---- circuito ----
  \draw (\xJ,0) to[R] (0.63,0) -- (0,0);
  \draw (0,-2.43) to[V, invert] (0,0);
  \node[ground] at (0,-2.43) {};
  \draw (\xJ,0) to[R] (\xJ,-2.43) node[ground]{};
  \draw (\xT,0) -- (\xE,0) -- (\xJ,0);
  \draw (\xT,0) to[R] (\xT,\yB);
  \draw (\xT,\yB) to[cI] (\xT,5.05) -- (\xT,\yC);
  \draw (\xT,\yB) -- (\xE,\yB) -- (\xBg,\yB) -- (\xBg,1.85) node[ground]{};
  \draw (\xT,\yC) -- (\xVo,\yC);
  \draw (\xRC,\yC) to[R] (\xRC,3.55) node[ground]{};
  \draw (\xRL,\yC) to[R] (\xRL,3.55) node[ground]{};
  \node[ground] at (\xVe,-2.225) {};
  \node[circ] at (\xJ,0) {};
  \node[circ] at (\xT,\yB) {};
  \node[circ] at (\xRC,\yC) {};
  \node[circ] at (\xRL,\yC) {};
  \node[ocirc] at (\xE,0) {};
  \node[ocirc] at (\xE,\yB) {};
  \node[ocirc] at (\xC,\yC) {};
  \node[ocirc] at (\xVo,\yC) {};

  % ---- rótulos dos componentes ----
  \node[anchor=east, inner sep=0] at (-0.574,-1.215) {$v_i$};
  \node[anchor=south, inner sep=0] at (1.74,0.40) {$R_S$};
  \node[anchor=west, inner sep=0] at (\xJ+0.343,-1.215) {$R_E$};
  \node[anchor=east, inner sep=0] at (\xT-0.343,1.225) {$r_e$};
  \node[anchor=west, inner sep=0] at (\xT+0.63,\yS) {$\alpha i_e$};
  \node[anchor=west, inner sep=0] at (\xRC+0.343,4.65) {$R_C$};
  \node[anchor=west, inner sep=0] at (\xRL+0.343,4.65) {$R_L$};

  % ---- terminais ----
  \node[term, anchor=base] at (\xC,\yC+0.264) {$C$};
  \node[term, anchor=base] at (\xE,\yB+0.264) {$B$};
  \node[term, anchor=base] at (\xE,0.264) {$E$};

  % ---- anotações ----
  \node[term, anchor=west, inner sep=0] at (\xVo+0.2,\yC) {$v_o = \alpha\, i_e\,(R_C \parallel R_L)$};
  \draw[term, line width=0.8pt, ->] (\xA,0.85) -- (\xA,1.6);
  \node[term, anchor=west, inner sep=0] at (\xA+0.2,1.225)
    {$i_e = \dfrac{v_e\rule[-5.7pt]{0pt}{5.7pt}}{r_e\rule{0pt}{10.3pt}}$};
  \draw[term, line width=0.8pt, ->] (\xRo,2.87) -- (\xRo,\yRo) -- (\xRC+0.30,\yRo);
  \node[term, anchor=north, inner sep=0] at (\xRo,2.73) {$R_o = R_C$};
  \draw[term, line width=0.8pt, ->] (\xRin,-3.2) -- (\xRin,-0.42) -- (\xJ-0.25,-0.42);
  \node[term, anchor=north, inner sep=0] at (\xRin,-3.32) {$R_{\mathrm{in}} = R_E \parallel r_e$};
  \node[term, anchor=base] at (\xVe,-0.912) {$+$};
  \node[term, anchor=base] at (\xVe,-1.424) {$v_e$};
  \node[term, anchor=base] at (\xVe,-2.011) {$-$};
\end{circuitikz}
\end{document}
